% ============================================================
% Table: IOC MR and MCL Violation Breakdown, 1985--2005 (Downstream CWSs)
% Purpose: Three-panel table describing violation codes and rule codes for
%          MR and MCL violations under IOC rules (nitrate 331, arsenic 332,
%          inorganic chemicals 333).
% Sample:  Strictly downstream CWSs (minehuc_downstream_of_mine=1, minehuc_mine=0)
% Source:  SDWA_VIOLATIONS_ENFORCEMENT.parquet + SDWA_REF_CODE_VALUES.csv
% N MR:    1,233 unique IOC MR violations, 1985--2005
% N MCL:   15 unique IOC MCL violations, 1985--2005
% ============================================================

\begin{table}[htbp]
\centering
\caption{IOC Violations by Violation Code and Regulatory Rule, 1985--2005 (Downstream CWSs)}
\label{tab:mr_mcl_violation_breakdown}
\small
\noindent\textbf{Panel A: MR violations by violation code} \\[4pt]
\begin{tabular}{clrr}
\hline\hline
\textbf{Code} & \textbf{Description} & \textbf{Count} & \textbf{Share (\%)} \\
\hline
\addlinespace[2pt]
\multicolumn{4}{l}{\textit{Routine and repeat monitoring failures}} \\
\addlinespace[2pt]
03 & Monitoring, Regular & 1,226 & 99.4 \\
04 & Monitoring, Check/Repeat/Confirmation & 7 & 0.6 \\
\addlinespace[2pt]
\hline
 & \textit{Total} & 1,233 & 100.0 \\
\hline\hline
\end{tabular}
\vspace{8pt}

\noindent\textbf{Panel B: MCL violations by violation code} \\[4pt]
\begin{tabular}{clrr}
\hline\hline
\textbf{Code} & \textbf{Description} & \textbf{Count} & \textbf{Share (\%)} \\
\hline
\addlinespace[2pt]
01 & Maximum Contaminant Level Violation, Single Sa.. & 10 & 66.7 \\
02 & Maximum Contaminant Level Violation, Average & 5 & 33.3 \\
\addlinespace[2pt]
\hline
 & \textit{Total} & 15 & 100.0 \\
\hline\hline
\end{tabular}
\vspace{8pt}

\noindent\textbf{Panel C: MR vs.~MCL violations by regulatory rule} \\[4pt]
\begin{tabular}{clrrrr}
\hline\hline
\textbf{Rule} & \textbf{Rule name} & \multicolumn{2}{c}{\textbf{MR}} & \multicolumn{2}{c}{\textbf{MCL}} \\
\cmidrule(lr){3-4}\cmidrule(lr){5-6}
\textbf{code} & & \textbf{Count} & \textbf{Share (\%)} & \textbf{Count} & \textbf{Share (\%)} \\
\hline
\addlinespace[2pt]
333 & Inorganic Chemicals & 844 & 68.5 & 11 & 73.3 \\
331 & Nitrate Rule & 291 & 23.6 & 2 & 13.3 \\
332 & Arsenic Rule & 98 & 7.9 & 2 & 13.3 \\
\addlinespace[2pt]
\hline
  & \textit{Total} & 1,233 & 100.0 & 15 & 100.0 \\
\hline\hline
\end{tabular}
\begin{minipage}{\linewidth}
\vspace{4pt}
\footnotesize
\textit{Notes:} Panel A: violation codes classify the type of monitoring-and-reporting failure; code 03 = Monitoring, Regular; 04 = Monitoring, Check/Repeat/Confirmation. Panel B: violation codes classify the type of MCL exceedance. Panel C: rule 331 = Nitrate; 332 = Arsenic; 333 = Inorganic Chemicals; MR = monitoring and reporting violation; MCL = maximum contaminant level violation; inorganic chemicals (333) encompasses nitrates and arsenic as sub-contaminants. All panels restricted to IOC rules (nitrate 331, arsenic 332, inorganic chemicals 333). Sample restricted to strictly downstream CWSs (minehuc\_downstream\_of\_mine\,=\,1 and minehuc\_mine\,=\,0), 1985--2005. Source: SDWA\_VIOLATIONS\_ENFORCEMENT.parquet, SDWA\_REF\_CODE\_VALUES.csv.
\end{minipage}
\end{table}
